\documentclass[11pt,a4paper]{article}
\usepackage{curso}
\begin{document}
\portada{Ejercicios --- Lección}{8.5}{K-FAC termodinámico}
\section{Ejercicios}
\begin{enumerate}[label=\textbf{E\arabic*.}]
\item (B) Razón entre el costo de inversión digital y termodinámico para $n=4096$, $\kappa=10$.
\item (B) Si la inversión es el 27\,\% del tiempo, ¿cuál es la aceleración máxima con el mismo número de pasos?
\item (B) ¿Por qué un cuantizador uniforme puede romper la semidefinición positiva?
\item (B) ¿Por qué los errores de salida pueden pesar más en entrenamiento que en un paso?
\item (I) Proponga un cuantizador que preserve la positividad (pista: aumente la diagonal en la cota del error de redondeo de cada fila).
\item (I) ¿Para qué $n$ el término $n^2\kappa^2$ es menor que $n^3$ con $\kappa=30$?
\item (A) Simule el entrenamiento de un perceptrón con K-FAC cuantizado y reproduzca la asimetría entrada/salida.
\end{enumerate}
\section{Soluciones}
\begin{solucion}
\textbf{E1.} $n^3/(n^2\kappa^2)=n/\kappa^2=40{,}96$.

\textbf{E2.} $1/(1-0{,}27)\approx1{,}37$.

\textbf{E3.} Redondear cada entrada puede volver negativo un valor propio pequeño: la matriz deja de ser SPD y la dinámica OU no converge.

\textbf{E4.} El error de salida entra en cada actualización y se acumula a lo largo de miles de pasos; un solo coseno no lo refleja.

\textbf{E5.} Sumar a cada diagonal la suma de los errores de redondeo máximos de su fila: por Gershgorin, la perturbación no hace negativo ningún valor propio.

\textbf{E6.} $n>\kappa^2=900$.

\textbf{E7.} Debe observarse que la salida requiere más bits cuando el cuantizador de entrada es conservador.
\end{solucion}
\end{document}
